\documentclass[../dayan_kernel_manual.tex]{subfiles}

\begin{document}

\section{Dayan API 参考}\label{sec:kernel-030}

本章中的函数由 \code{dayan_sys.h} 声明，是算子 kernel 可直接链接的 Dayan API。
所有返回 \code{int} 的 resource 接口以 0 表示成功、负值表示失败；公开头文件当前
没有导出 errno 名称，算子应使用 \code{status < 0} 判断失败，不应依赖具体负数值。

\subsection{协作式让出}

\begin{dayantable}{调度接口}{tab:kernel-030}{X[27,l] X[73,l]}
  符号 & 语义 \\
  \code{dayan_yield} & 将当前 proc 重新放到本 hart ready queue 尾部并切回 scheduler；再次被选择后从调用点继续。 \\
\end{dayantable}

\code{dayan_yield} 不保证其它特定 proc 一定执行，也不是内存屏障或跨核同步原语。
如果当前 proc 是唯一 ready 的普通任务，仍可能很快再次运行。

\subsection{张量单元等待}

\begin{dayantable}{Produce/Bar 等待接口}{tab:kernel-031}{X[27,l] X[18,l] X[55,l]}
  符号 & 单元 & 语义 \\
  \code{dayan_wait_lsu} & LSU & 等待调用前已经发射到 LSU 的异步工作完成。 \\
  \code{dayan_wait_dma} & LSU/DMA & \code{dayan_wait_lsu} 的兼容别名；DMA 由 LSU 路径实现。 \\
  \code{dayan_wait_mat} & MAT & 等待调用前已经发射到矩阵执行单元的工作完成。 \\
  \code{dayan_wait_vec} & VEC & 等待调用前已经发射到向量执行单元的工作完成。 \\
  \code{dayan_wait_macro} & MACRO & 等待调用前已经发射到宏指令执行单元的工作完成。 \\
\end{dayantable}

这些函数把 \code{tp} 的当前值直接作为 Produce token 写入硬件队列，然后协作式切出。
返回表示对应 token 已从全局完成队列取出，不表示其它单元或调用之后发射的工作完成。

\begin{dayancode}
issue_vec_operations();
dayan_wait_vec();
consume_vec_results();
\end{dayancode}

不得以短循环无界发出 wait：全局 Produce 队列保证容量不少于 32 项，每单元队列保证
容量不少于 2 项，但超过该上限的软件行为可能阻塞硬件。

\subsection{单核 resource}

\code{dayan_resource_t} 只允许同一 hart 上的 proc 共享。其操作不使用原子指令；
Runtime 利用非抢占执行和对象内嵌的事件链表完成等待与唤醒。

\begin{dayantable}{单核 resource 接口}{tab:kernel-032}{X[34,l] X[66,l]}
  符号 & 语义 \\
  \code{dayan_resource_init} & 初始化为 FREE，计数清零并初始化三条等待链表。对象首次使用前调用一次。 \\
  \code{dayan_resource_alloc} & 仅在 FREE 时设置 producer/consumer 数量；consumer 必须非零，producer 可以为零。 \\
  \code{dayan_resource_produce} & 在 BUSY 状态减少 producer 数量；最后一个 producer 把状态改为 READY 并唤醒全部 ready waiter。 \\
  \code{dayan_resource_consume} & 在 READY 状态减少 consumer 数量；最后一个 consumer 把状态改为 FREE 并唤醒全部 free waiter。 \\
  \code{dayan_resource_wait_ready} & 状态不是 READY 时，把当前 proc 挂入该对象的 ready 等待链表；可能多次睡眠和复查。 \\
  \code{dayan_resource_wait_free} & 状态不是 FREE 时，把当前 proc 挂入该对象的 free 等待链表；可能多次睡眠和复查。 \\
\end{dayantable}

对象可以位于当前 hart 的 proc 栈或私有堆。初始化后，只要可能存在等待者，就不得复制、
移动、释放或再次初始化对象；其 \code{list_head} 字段保存的是实际 proc 节点链接。
字段公开是为了固定 C ABI，不表示算子可以直接修改状态、计数、锁或链表。

一个典型使用模式如下：

\begin{dayancode}
int prepare_resource(dayan_resource_t *resource)
{
    int status = dayan_resource_init(resource);
    if (status < 0)
        return status;
    return dayan_resource_alloc(resource, 2, 1);
}

void producer_done(dayan_resource_t *resource)
{
    if (dayan_resource_produce(resource) < 0)
        return;
}

void consumer(dayan_resource_t *resource)
{
    if (dayan_resource_wait_ready(resource) == 0)
        (void)dayan_resource_consume(resource);
}
\end{dayancode}

\subsection{跨核 shared resource}

\code{dayan_shared_resource_t} 用于多个 hart 通过 coherent window 共享状态。对象必须
放在当前 proc 的 \code{shared} allocation 中，并由一个明确的初始化者在发布指针前
完成初始化。

\begin{dayantable}{跨核 resource 接口}{tab:kernel-033}{X[38,l] X[62,l]}
  符号 & 语义 \\
  \code{dayan_shared_resource_init} & 校验对象位于 \code{0x00400000--0x007fffff} 且不低于当前 \code{shared}，原子初始化为 FREE。 \\
  \code{dayan_shared_resource_alloc} & 在原子锁内由 FREE 建立 producer/consumer 计数。 \\
  \code{dayan_shared_resource_produce} & 在原子锁内减少 producer 计数，最后一次操作以 release 语义发布 READY。 \\
  \code{dayan_shared_resource_consume} & 在原子锁内减少 consumer 计数，最后一次操作以 release 语义发布 FREE。 \\
  \code{dayan_shared_resource_wait_ready} & 以 acquire load 检查 READY；未满足时进入 Runtime poll wait。 \\
  \code{dayan_shared_resource_wait_free} & 以 acquire load 检查 FREE；未满足时进入 Runtime poll wait。 \\
\end{dayantable}

\begin{dayancode}
struct shared_layout {
    dayan_shared_resource_t resource;
    uint32_t payload[64];
};

struct shared_layout *layout = (struct shared_layout *)shared;

/* Exactly one participant, before publishing the object: */
int status = dayan_shared_resource_init(&layout->resource);
\end{dayancode}

初始化函数只能验证硬件窗口边界，无法知道 Host 给当前 proc 分配的精确长度。Host
必须保证整个对象位于 allocation 内，并保证多个参与 hart 看到的是同一个 coherent
物理对象。

shared resource 不内嵌 proc 等待链表，因为普通 NPU 缓存中的 proc 指针属于单个
Runtime 实例，不能作为跨核一致性对象的一部分。等待使用 scheduler poll，状态变化
不需要也不能从另一个 hart 直接操作本核 ready queue。

\subsection{错误与所有权规则}

\begin{itemize}
  \item 空指针、consumer 数量为 0 或 shared 地址非法会失败；
  \item 在错误状态调用 alloc/produce/consume 会失败，不会自动修复状态机；
  \item resource 函数不会输出错误信息，也不会终止 proc；
  \item 同一个 \code{dayan_resource_t} 不得跨 hart 使用；
  \item 同一个 \code{dayan_shared_resource_t} 不得由多个 hart 并发初始化；
  \item Host 回收 shared allocation 前，必须确保所有使用该对象的 proc 已结束且没有
        hart 仍在轮询。
\end{itemize}

\end{document}
